\documentclass[11pt]{article}
\usepackage{kotex,amsmath,amssymb,booktabs}
\usepackage[margin=1in]{geometry}
\usepackage[hidelinks]{hyperref}
\title{JLZ v13 realized writer: pilot와 2,000-edit 순차 실험 설계}
\author{실행 준비 명세}
\date{2026년 10월 5일}
\begin{document}
\maketitle

\paragraph{실행 상태.}
본 문서는 다음 실험을 준비한다. 현재 AlphaEdit/ridge job을 취소하거나 수정하지
않는다. 실제 GPU pilot 또는 본 실험을 제출했다는 의미가 아니다.
원하는 관측은 native subject에서의 계획 실현, 실제 배분, 정확도와 순차 보존이다.
기술 qualification과 과학적 성과를 구분하고 PS/NS/배분 다양성으로 실행을 거르지 않는다.

\section{질문과 동결 조건}
질문은 (1) equality/range-LS writer가 soft ridge보다 local 목표를 잘 실현하는가,
(2) direct 정책이 계획 외 보정량을 줄이는가,
(3) 평균 제약보다 native context 제약이 실제 주입 위치의 오차를 줄이는가,
(4) 그 결과가 ACC, paraphrase, locality와 순차 retention에 어떤 영향을 주는가다.
실현이 좋아졌다고 task/locality 성능이 자동으로 좋아졌다고 판정하지 않는다.

v12의 전체 eligible 층, native rewrite/KL 문장, tokenization, teacher,
손실과 계수, .75 예산, EfficiencyAdam, 요청별 stop/terminal 규칙을 동결한다.
주 profile은 동일 Llama3-8B-Instruct/CounterFact first2000 순서, L4--L8,
FP32 모델, FP64 writer geometry, TF32 off, 같은 C0와 초기 H=0이다.
정확한 model/data/context/source/runtime SHA는 실제 제출 전 manifest에 결속한다.
현재 profile의 native KL row까지 제약하는 CD와 rewrite 평균만 제약하는 MD를
명시하며, history는 두 경우 모두 rewrite-only native 평균으로 유지한다.

\section{Q0: CPU 수학 참조}
작은 행렬로 full-rank 등식/최소 에너지, compatible duplicate, 서로 다른 목표의
duplicate, zero-target 열, context mismatch, metric whitening, numerical rank,
FP32 applied update, variable owner/context 수와 마지막 batch를 검증한다.
KL의 목표는 같은 owner D이며 history 평균에서 KL은 제외한다.
결과는 \texttt{reference/audit.json}에 기록한다. CPU PASS는 실제 모델이나
속도 qualification을 대신하지 않는다.

\section{Q1: 작은 실제 모델 qualification}
동일 profile의 first4 요청을 BS2로 두 번 순차 적용한다. 각 MD/CD chain은
자기 W/H에서 두 번째 계획을 만든다. 목적은 native 행 구성, actual upper key
재측정, write orientation, solver residual, 실제 FP32 작용, H exactly-once,
실패 rollback 및 observer 비변이를 검증하는 것이다.
실제 B100의 rank/시간/품질은 이 작은 검사로 대신하지 않는다.
0 목표, 중복 key 충돌과 near-singular 문제는 CPU에서 실패/근사를 의도적으로
검사하며 실제 benchmark label/목표를 변조해 qualification하지 않는다.

\section{P1: first100의 같은 계획으로 writer를 비교}
Cold W0/H0에서 v12 계획을 한 번만 fit한다. 동일 $u,D,z^v$, entry teacher,
context/token metadata를 고정하고 다음 5개 branch를 실행한다.

\begin{center}
\begin{tabular}{lll}
\toprule
Branch & Writer와 key & 우변 정책\\
\midrule
RT & 기존 MEMIT-H soft ridge, native mean & $z_l^v-h_l^{cur}$ tracking\\
RD & 같은 soft ridge, native mean & $D_l$ direct\\
MT & 최소 A-에너지 equality/range-LS, native mean & $z_l^v-h_l^{cur}$ tracking\\
MD & 최소 A-에너지 equality/range-LS, native mean & $D_l$ direct\\
CD & 최소 A-에너지 equality/range-LS, fit injection rows & owner별 $D_l$ direct\\
\bottomrule
\end{tabular}
\end{center}

RT/RD/MT/MD의 2$\times$2는 동일 보존 metric에서 writer와 residual 정책을
분리한다. MD/CD는 제약 위치의 효과를 분리한다. AlphaEdit을 이 2$\times$2의
soft arm으로 쓰면 metric까지 바뀌므로 사용하지 않는다. 이전 AlphaEdit B1은
참고 자료로 재사용하되 신규 branch인 것처럼 집계하지 않는다.

각 branch는 동일 W/H/RNG에서 시작하고, 하층 write 이후 자신의 상층 key를
재계산한다. 공통 계획만 공유하며 다른 branch의 K/hidden은 공유하지 않는다.
모든 층 write와 native final H append 후 동일 R/P/N/ACC를 평가하고 RAM에서
W/H/RNG/cache를 원복한다. 원복 hash와 plan hash를 확인하며 추가 fit은 0이다.
모든 control branch의 비용과 main 비용은 따로 기록한다.

성능 결과를 보고 budget, rank threshold, coefficient, writer를 바꾸지 않는다.
미리 정한 main은 아래 MD와 CD 두 arm이다. 기술적으로 검증된 finite range-LS는
rank 부족을 기록하고 계속한다. Numerical realization failure 등 기술 오류가 나면
그 branch를 미검증 상태로 남기고 원인을 고친 새 source로 재qualification한다.
실패 branch를 조용히 ridge로 치환하거나 성공 요청만 다음 단계에 넘기지 않는다.

\section{Main: MD와 CD 각각 2,000 edit}
두 arm 모두 독립된 cold W0/H0에서 first2000을 시작한다. 기본 논리 batch는
BS100이며 총 20회다. MD는 평균 exact/LS direct, CD는 native-context exact/LS
direct다. Main에서 RT/RD/MT의 추가 2k fit을 하지 않는다.
P1의 plan을 이후 다른 entry 상태에서 재사용하지 않는다. B2부터 각 arm은 자기
W/H를 입력으로 새 공동 계획을 fit한다. 따라서 장기 결과는 각 writer를 사용한
전체 순차 방법의 효과이고, 같은 계획 단독 효과는 P1에서 판정한다.

기존 v12 ridge first2k는 source/runtime/input/order/evaluator manifest가 맞을 때
기존 실행 대조군으로 재사용한다. Native MEMIT-H, AlphaEdit, BLUE 등은 일치하지
않는 runtime/입력 조건이 있으면 historical reference로 표시한다. 새 baseline
fit을 이번 설계에 묵시적으로 추가하지 않는다.

\subsection{Batch, 모델, benchmark 독립성}
RunSpec은 총 편집 수 $N$, 논리 $B$ 또는 명시적 schedule을 받는다. 모든 구조는
실제 $B_t$, 각 요청 context 수와 layer width에서 만든다. BS128이면
$128\times15+80=2000$의 16회이며 마지막 요청을 버리지 않는다.
N과 B는 실행 profile이며 method 상수가 아니다. 두 main arm은 같은 schedule을 쓴다.
Fit request 묶음, capture row 묶음, evaluation row 묶음은 독립적으로 메모리에
맞춘다. Writer는 전체 logical batch의 제약을 함께 풀며 microbatch write로 쪼개지 않는다.
모델별 module, additive mapping, key dependency와 benchmark별 row/weight/metric은
adapter qualification 대상이다. 논리 B 변경이 결과 불변을 의미하지는 않는다.

\section{관측과 판정}
매 commit에서 current cohort를 평가하고, 누적 edit 수 500/1000/1500/2000에서
all-seen 평가를 수행한다. W0에서 같은 first2000의 reference를 한 번 계산한다.
기본 BS100의 R/P/N 분모는 각 시점 실제 평가 row로 검산한다. 다른 benchmark에서는
2P/10N 상수를 쓰지 않는다. Milestone이 batch 경계와 다르면 다음 완료 commit의
실제 edit 수를 표시하며, 정확한 중간 수가 필요하면 schedule에 미리 경계를 넣는다.

주 평가 항목은 RS/PS/NS, R/P desired-new 및 N desired-true의 TF strict ACC,
token-micro와 prompt-macro ACC, new/true NLL과 margin이다. TF를 자유생성 정확도로
부르지 않는다. R/P/N 선호율의 harmonic score는 보조 지표다.
요청 identity로 편집 직후와 W5/W10/W15/W20을 연결하고, preference와 strict ACC
각각의 lost/gained/retained를 기록한다. Active/superseded는 별도 층화하고 전체
분모에서 임의로 제거하지 않는다. 이 평가 값은 optimizer나 writer 선택에 입력하지 않는다.

\paragraph{실현과 배분.}
각 층/요청에서 ideal, effective-weight, 실제 forward의 직접 작용을 구분한다.
Native mean/canonical/rewrite context/KL을 나눠 norm ratio, 방향 투영비, cosine,
목표 오차를 기록한다. 0 목표는 비율 대신 anchor 정규화 leakage로 기록한다.
Plan share와 mean/canonical/context 직접 share, 요청별 share L1 오차,
상속 gap, final virtual-target gap과 실제 entry 대비 순변화를 따로 남긴다.
최종 gap은 per-request 벡터로 직접 계산하고 scalar median의 곱으로 대체하지 않는다.

\paragraph{Feasibility와 비용.}
각 solve의 shape, numerical rank, singular cutoff/condition, 버린 목표 성분,
compatible 여부, ideal projected-target residual, applied-target residual,
\(Q=\operatorname{tr}(UAU^\top)\), update norm을 남긴다.
이상적 Q는 retained coefficient의 Frobenius norm으로 저비용 계산한다.
실제 FP32 $Q_{eff}$는 P1 전체 및 main 100/500/1000/1500/2000 milestone에서
별도로 계산하고 나머지는 \texttt{NOT\_MEASURED}로 표시한다. 이상값을 실제 에너지로
재표기하지 않는다. 성능이나 에너지의 임의 임계값은 채택 gate가 아니다.

\paragraph{해석.}
국소 실현 오차가 감소했어도 PS/ACC/NS/retention이 개선되지 않으면 성공으로
요약하지 않는다. CD가 MD보다 실현을 개선하지만 NS를 악화시키면 stronger context
제약의 tradeoff로 보고한다. 두 arm의 final strength가 같다고 가정하지 않고 실제
subject/context 변위와 logits/NLL을 함께 보고한다. 이 실험만으로 배분 정책의
독립적 신규성이나 최적성을 증명하지 않는다.

\section{실행 비용과 인계}
현재 profile은 native row 700개가 이미 forward되므로 KL key를 저장하는 것은
추가 capture forward를 요구하지 않는다. 제약 열이 MD의100에서 CD의700으로 늘어
triangular apply 및 rank decomposition 비용은 증가한다. 전체 시간이 7배라는 뜻은 아니다.
기존 ridge2k의 fit 시간 약7789초는 workload 참고치이며 새 writer의 ETA가 아니다.
P1에서 capture/factor/SVD/apply/history/observer, GPU peak/CPU RSS를 분리 측정한다.
Main ETA는 해당 실측과 누적 평가 문항 수로 계산하고 초기 수치를 확정 ETA로 제시하지 않는다.

동시에 main 두 개를 올려 경쟁시키지 않고 가용 자원에서 직렬 실행하는 것을 기본으로
한다. 제안 owner/server는 직전 작업의 SH3/server3이며 실제 자원과 source lock은
후속 실행 지시 시 확인한다. 기존 job은 건드리지 않는다. Persistent model checkpoint를
새로 추가하지 않으며, CPU scalar/raw metrics와 source/manifest를 저장한다.
계획/branch 비교의 tensor snapshot은 RAM transaction으로 제한한다.

구현 완료 후 source/config/model/data/context/row order/metric SHA, qualification
receipt, branch 복원 결과, main config 두 개를 결속한 패키지를 GH/SH3에 인계한다.
이번 문서 작성 자체는 해당 전달이나 GPU 제출을 수행한 것이 아니다.
\end{document}
